\section{Finite calculations}\label{app:finite-calculations}

This appendix gives the finite quantities used above.
Each terminating decimal denotes an exact rational number.
All decimal interval endpoints are rounded outwards; decimal approximations
in explanatory prose are not additional certified bounds.
The comparisons with prime sums occur in the main text.

\subsection{Series bounds}

For \(|z|<1\),
\begin{equation}\label{eq:app-log-tail}
 \left|\operatorname{arctanh}z-
 \sum_{k=0}^{m-1}\frac{z^{2k+1}}{2k+1}\right|
 \le\frac{|z|^{2m+1}}{(2m+1)(1-z^2)}.
\end{equation}
Write any positive rational as \(t=2^a v\),
where \(2/3\le v<4/3\).
Then use
\begin{equation}\label{eq:app-log-reduction}
 \log t=a\log2+2\operatorname{arctanh}\frac{v-1}{v+1},
 \qquad \log2=2\operatorname{arctanh}(1/3).
\end{equation}
Also, for \(0\le z<1\),
\begin{equation}\label{eq:app-dilog-tail}
 0\le\operatorname{Li}_2(z)-\sum_{n=1}^m\frac{z^n}{n^2}
 \le\frac{z^{m+1}}{(m+1)^2(1-z)}.
\end{equation}
These inequalities enclose the elementary functions below.
Integer comparisons then give the printed rational bounds.

\subsection{Four profile differences}\label{subsec:app-h6-four-profiles}

Put
\begin{gather}
 \alpha=\frac{100}{1149},\quad \beta=\frac{50}{309},\quad
 \mu_* =\frac{25}{128},\quad \lambda=\frac{93}{400},\quad
 h_G=\frac9{5000},\quad K=\frac{57}{64}-\frac1{500},\notag\\
 \delta_h=\frac {h_G}\alpha,\qquad C=K-\frac{11}{5}\alpha,\qquad
 B(x)=\min\{\beta,\max\{\alpha,C-5x/2\}\}.
 \label{eq:app-h6-profile-parameters}
\end{gather}
Define the elementary functions
\begin{align}
 \widehat f(s)&=\frac{\log(s-1)}s &&(2\le s\le4),\notag\\
 \widehat F(s)&=\frac1s &&(1<s\le3),\label{eq:app-h6-profile-functions}\\
 s\widehat F(s)&=1+\log(s-2)\log(s-1)
 -\operatorname{Li}_2\!\left(\frac{s-2}{s-1}\right)\\
 &\quad-\frac{\log^2(s-1)}2+\operatorname{Li}_2(1/2)
 +\frac{\log^22}{2} &&(3<s<5),\notag
\end{align}
The two sieve normalizations satisfy
\[
 \widehat F=F/(2e^\gamma),\qquad \widehat f=f/(2e^\gamma).
\]
The raw multiplier $2/(e^\gamma\alpha)$, a signed row coefficient
$\pm1/4$, and $F=2e^\gamma\widehat F$ or
$f=2e^\gamma\widehat f$ combine to the displayed $1/\alpha$ factor.
The directed integrals below use these hatted functions. No further $1/4$ or
$2e^\gamma$ factor is applied.
Also $\operatorname{Li}_2(t)=\sum_{n\ge1}t^n/n^2$ for $0\le t<1$.
Write
\[
 D_F(s)=\widehat F(s)-\widehat F(s+\delta_h),\qquad
 D_f(s)=\widehat f(s+\delta_h)-\widehat f(s).
\]
For the three intervals in Table~\ref{tab:app-h6-profile-bounds}, put
\begin{align}
 \Delta_j&=\frac1\alpha\int_{I_j}
 D_F\!\left(\frac{K-5x/2}{\alpha}\right)\frac{dx}{x}
 \quad(j=4,5,7),\label{eq:app-h6-profile-negative}\\
 \Delta_{10}&=\frac1\alpha\int_{\mu_*}^\lambda\int_\alpha^{B(x)}
 D_f\!\left(\frac{K-5x/2-y}{\alpha}\right)
 \frac{dy\,dx}{xy}.
 \label{eq:app-h6-profile-positive}
\end{align}
These $\Delta_j$ are those of
equation~(5.22).
Closed or half-open interval endpoints have measure zero.
The curved domain is that of \cite[(2.38)--(2.40)]{PaperII}.

\begin{proposition}\label{prop:app-h6-four-profiles}
A directed interval computation certifies the following outward enclosure of the
four elementary integrals:
\begin{equation}\label{eq:app-h6-profile-total}
 \frac{572004}{10^8}<\Delta_4+\Delta_5+\Delta_7+\Delta_{10}
 <\frac{606486}{10^8}.
\end{equation}
In particular, their sum minus $10^{-20}$ exceeds $5719/10^6$.
Also,
\begin{equation}\label{eq:app-h6-profile-strip}
 \frac{\log3}{4\alpha^2}
 \left(\frac52\lambda+\beta\right)\frac{\lambda-\mu_*}{\mu_*}
 <\frac{21137254516197687}{4120000000000000}
 <\frac{5137}{1000}<\frac{21}{4}<\frac{179}{25}.
\end{equation}
\end{proposition}

\begin{table}[ht]
\centering
\caption{Outward decimal bounds from the directed computation. Every
endpoint is an exact terminating rational.}
\label{tab:app-h6-profile-bounds}
\begin{tabular}{cccc}
\hline
$j$ & $I_j$ & lower bound & upper bound\\
\hline
$4$ & $[25/128,1/4]$ & $0.00162068$ & $0.00169539$\\
$5$ & $[1/4,7109/28000]$ & $0.00040327$ & $0.00041325$\\
$7$ & $[25/128,93/400]$ & $0.00056463$ & $0.00059051$\\
$10$ & curved domain & $0.00313145$ & $0.00336572$\\
\hline
sum & & $0.00572003$ & $0.00606487$\\
\hline
\end{tabular}
\end{table}

\begin{proof}
The profiles are evaluated with directed interval arithmetic. An independent
verifier rechecks the directed output. Put
$h_G=9/5000$, $p(s)=-\widehat F'(s)$, and let $d$ be the right-branch
derivative profile for $\widehat f$, so $d(s)=\widehat f'(s)$ for $2<s<4$
and $d(2)=1/2$ by continuous extension.
The three one-dimensional rows each use 2048 equal cells:
\[
 x_k=a_j+\frac{k(b_j-a_j)}{2048}\qquad(0\le k\le2048).
\]
For a cell $[x_0,x_1]=[x_k,x_{k+1}]$, define
\[
 s_-=\frac{K-5x_1/2}{\alpha},\qquad
 s_+=\frac{K-5x_0/2+h_G}{\alpha}.
\]
Since $D_F(s)=\int_0^{h_G/\alpha}p(s+t)\,dt$, the code's lower and upper
cell contributions are exactly
\begin{equation}\label{eq:app-h6-profile-negative-cells}
 L=\frac{h_G}{\alpha^2}\frac{x_1-x_0}{x_1}p(s_+),\qquad
 U=\frac{h_G}{\alpha^2}\frac{x_1-x_0}{x_0}p(s_-).
\end{equation}
These hold because $s_-\le (K-5x/2)/\alpha+t\le s_+$ for every
$x\in[x_0,x_1]$ and $0\le t\le h_G/\alpha$; also $x_0\le x\le x_1$.

For the fourth integral, split at
\[
 x_B=\frac{506848823}{2366940000}.
\]
On each piece $[a,b]$ of $[\mu_*,x_B]\cup[x_B,\lambda]$, the S10 routine uses
$x_i=a+i(b-a)/96$ and $z_j=j/96$ for $0\le i,j\le96$.
For a cell $[x_i,x_{i+1}]\times[z_j,z_{j+1}]$, with $0\le i,j<96$, put
\[
 b_0=B(x_i),\quad b_1=B(x_{i+1}),\quad
 y_-=\alpha+(b_1-\alpha)z_j,\quad
 y_+=\alpha+(b_0-\alpha)z_{j+1},
\]
\[
 q_-=\frac{K-5x_{i+1}/2-y_+}{\alpha},\qquad
 q_+=\frac{K-5x_i/2-y_-+h_G}{\alpha}.
\]
The substitution is $y=\alpha+(B(x)-\alpha)z$. Its executed lower and
upper contributions are
\begin{equation}\label{eq:app-h6-profile-positive-cells}
 L=\frac{h_G}{\alpha^2}(x_{i+1}-x_i)(z_{j+1}-z_j)
       \frac{b_1-\alpha}{x_{i+1}y_+}d(q_+),\qquad
 U=\frac{h_G}{\alpha^2}(x_{i+1}-x_i)(z_{j+1}-z_j)
       \frac{b_0-\alpha}{x_iy_-}d(q_-).
\end{equation}
For every point in this cell and $0\le t\le h_G/\alpha$,
$q_-\le (K-5x/2-y)/\alpha+t\le q_+$, and
\[
 \frac{b_1-\alpha}{x_{i+1}y_+}
 \le\frac{B(x)-\alpha}{xy}
 \le\frac{b_0-\alpha}{x_iy_-}.
\]
Hence the two positive factors are bounded separately, as in the code.
The two rectangular grids contain $2\cdot96^2$ cells. The total
executed cell count is
\[
 3\cdot2048+2\cdot96\cdot96=24576.
\]

The endpoint bounds use the following branch formulas. For $2<s<4$,
\[
 d(s)=\frac{s/(s-1)-\log(s-1)}{s^2}>0,\qquad
 d'(s)=-\frac1{s(s-1)^2}
       -\frac{2\{s/(s-1)-\log(s-1)\}}{s^3}<0.
\]
They extend continuously to $s=2$ with $d(2)=1/2$ and right derivative
$d'_+(2)=-1$. Thus the branch profile is positive and decreasing on $[2,4]$.
For positivity: the numerator decreases, and at $s=4$ it is
$4/3-\log3>0$, since
$e^{4/3}>1+4/3+(4/3)^2/2+(4/3)^3/6>3$. For $s>3$, put
$A(s)=\widehat F(s)-\widehat f(s-1)$. Then
\[
 (sA(s))'=-s\,\widehat f'(s-1)<0,\qquad
 p(s)=A(s)/s,\qquad
 p'(s)=-2A(s)/s^2-\widehat f'(s-1)/s.
\]
At $s=3$ the two pieces give $p(3)=1/9$, so $p$ is continuous
there; the displayed differential identities are used on $s>3$.
The exact directed evaluation at
$s_{\max}=73930107/16000000$ gives
$A(s_{\max})>1597/100000>0$. Hence $A(s)>0$ throughout the required
range $3<s\le s_{\max}$, and there $p>0$ and $p'<0$. For $s\le3$,
$p(s)=1/s^2$ is also positive and decreasing. The S10 mixed-corner range is
\[
 \frac{17359415273}{7910400000}\le q_-\le q_+
 \le\frac{57930107}{16000000}<4,
 \qquad \frac{17359415273}{7910400000}>\frac{219}{100}.
\]
These inequalities justify using the decreasing positive $d$ in
\eqref{eq:app-h6-profile-positive-cells}. The lower bound refers only to the
executed mixed-corner interval.

The code evaluates the formulas by rational geometry and 256-bit Arb
ball arithmetic. In \texttt{profile\_direct.py}, logarithms use Arb's
\texttt{log()} operation. Put $w=s-2$. The dilogarithm uses Arb's
\texttt{polylog(2)} operation at the negative argument $-w$, together with
$\pi^2/12$. These are the code's evaluations of
\eqref{eq:app-h6-profile-functions}. By the Landen identity
\[
 \operatorname{Li}_2(-w)
 =-\operatorname{Li}_2\!\left(\frac{w}{1+w}\right)
  -\frac12\log^2(1+w),\qquad
 \operatorname{Li}_2(1/2)+\frac12\log^22=\frac{\pi^2}{12}
\]
the Arb expression equals the displayed formula. Arb supplies
outward ball enclosures. The code does not use the atanh and
dilogarithm tail inequalities above. The exact directed rational
endpoints are written to the result record. The terminating decimals in the
table are outward projections of each row's directed endpoint. Their displayed
column sums are $0.00572003$ and $0.00606487$. The sharper bound
\eqref{eq:app-h6-profile-total} adds the unrounded rational endpoints
in RESULT.json before decimal projection. Summing before rounding proves the stated
integer comparisons. Finally, 64 terms of
$\log3=2\sum_{n\ge0}2^{-(2n+1)}/(2n+1)$ and its geometric remainder give
$\log3<10987/10000$. This proves \eqref{eq:app-h6-profile-strip}.
\end{proof}

Proposition~\ref{prop:app-h6-four-profiles} concerns elementary integrals and
one strip coefficient. Its application to the prime sums requires the
source comparisons, uniform errors and a single charge of the strip.


\subsection{The shifted three-variable integral}\label{subsec:app-m5}

Here \(F,f\) are the linear-sieve functions.
Their normalization is
\begin{equation}\label{eq:app-delay-functions}
\begin{gathered}
 F(s)=2e^\gamma/s\quad(1\le s\le3),\qquad f(s)=0\quad(s\le2),\\
 (sF(s))'=f(s-1)\quad(s>3),\qquad(sf(s))'=F(s-1)\quad(s>2),\\
 0\le f\le1\le F,\qquad f\text{ is nondecreasing}.
\end{gathered}
\end{equation}
These classical properties are the inputs to this calculation.
Put \(u=1/500\), \(\mu_*=25/128\) and \(h_G=9/5000\).
Define
\begin{equation}\label{eq:app-six-branch}
\begin{aligned}
 v_1(s)&=\min\left(\frac{1+s}{2},\frac{32-34s}{39},
                         \frac{57}{64}-\frac{3s}{2}\right),&
 v(s)&=v_1(s)-u,\\
 \psi(x,y)&=\begin{cases}x,&x\ge y,\\-\infty,&x<y,\end{cases}\\
 \Theta_*(x,y,z)&=\max\{v(x),v(y),v(x+y),v(x+y+z),\\
 &\hspace{2.7em}\psi(v(x+z),x+2y+z),
               \psi(v(y+z),2x+y+z)\}.
\end{aligned}
\end{equation}
The unguarded branches include \(v(x)>0\) for \(u\le x\le\mu_*\).
Thus replacing the inactive value \(-\infty\) by \(0\)
leaves this maximum unchanged.
For \(d=0,h_G\), put
\begin{equation}\label{eq:app-G-functional}
\begin{aligned}
 J_d&=\int_{u\le z\le y\le x\le\mu_*}
 \frac{1-f((\Theta_*(x,y,z)+d-x-y-z)/z)}{xyz^2}
 \,dx\,dy\,dz,\\
 H_{\log}&=\log(\mu_*/u),\\
 R_{\mathrm{up}}&=500(F(30)-1)+500H_{\log}(1-f(30))
                   +250H_{\log}^2(F(30)-1),\\
 \mathcal G_d&=1/\mu_*+R_{\mathrm{up}}+J_d.
\end{aligned}
\end{equation}
The unshifted $d=0$ profile and functional come from
\cite[(4.70)--(4.75)]{PaperII}. The parameter $d$ and the shifted functional
$\mathcal G_d$ are defined in this paper. The retained-side argument uses the
instance $d=h_G$, constructed in
\eqref{eq:retB-G-use}. Paper~II's (5.4)--(5.5) are not the source of this
shifted functional.

\begin{proposition}\label{prop:app-m5-cap}
The shifted scalar satisfies
\begin{equation}\label{eq:app-m5-cap}
 0\le\mathcal G_{h_G}\le U_h<\frac{6073714261}{10^9}.
\end{equation}
Let \(F_h^+\ge F(30)\), \(f_h^-\le f(30)\),
and \(H_h^+\ge H_{\log}\) be the directed table endpoints.
Define
\begin{equation}\label{eq:app-directed-tail}
\begin{aligned}
 \mathcal R(X,Y,Z)&=500(X-1)+500Z(1-Y)
                             +250Z^2(X-1),\\
 R_h&=\mathcal R(F_h^+,f_h^-,H_h^+),\\
 U_h&=1/\mu_*+R_h+J_h^+.
\end{aligned}
\end{equation}
The construction has \(F_h^+\ge1\), \(f_h^-\le1\).
Thus \(R_{\mathrm{up}}\le R_h\). This directed endpoint comparison is the
monotonicity step of \cite[\S4.6, (4.73)]{PaperII}.
Table~\ref{tab:app-m5} encloses these finite upper bounds.
The constants $K_{\mathrm{hi}}$ and $L_{G1}$ are printed exactly in
\eqref{eq:retB-exact-Khi} and \eqref{eq:retB-exact-LG1}, with their
Paper~II locators.
The normalization identities are
\begin{equation}\label{eq:app-m5-normalization}
\begin{aligned}
 M_{5,\mathrm{internal}}
 &=-K_{\mathrm{hi}}\frac{6073714261}{10^9}-L_{G1},\\
 M_{5,\mathrm{internal}}+
 K_{\mathrm{hi}}\left(\frac{6073714261}{10^9}-U_h\right)
 &=-K_{\mathrm{hi}}U_h-L_{G1}.
\end{aligned}
\end{equation}
\end{proposition}

\begin{table}[ht]
\centering\small
\caption{Directed bounds for the shifted profile with $n_{\mathrm{mesh}}=512$.}
\label{tab:app-m5}
\begin{tabular}{@{}lll@{}}
\hline
Quantity & lower endpoint & upper endpoint\\
\hline
\(R_h\) & \(0.000004373236955\) & \(0.000004373236956\)\\
\(J_h^+\) & \(0.953709886776891\)
 & \(0.953709886776892\)\\
\(U_h\) & \(6.073714260013847\)
 & \(6.073714260013848\)\\
\(6073714261/10^9-U_h\)
 & \(0.000000000986152\)
 & \(0.000000000986153\)\\
\hline
\end{tabular}
\end{table}

\begin{proof}
For the outer partition put \(n=512\), \(x_0=u\), \(x_n=\mu_*\), and
\begin{equation}\label{eq:app-m5-mesh}
 x_i=2^{-48}\left\lfloor2^{48}
       \left(u+(\mu_*-u)\sqrt{i/n}\right)\right\rfloor
       \quad(1\le i<n).
\end{equation}
It is strictly increasing.
For the inner integral write
\begin{equation}\label{eq:app-inner}
 g_d(x,y)=\int_u^y
 \frac{1-f((\Theta_*(x,y,z)+d-x-y-z)/z)}{z^2}\,dz.
\end{equation}
Each branch residual decreases in \(x,y\).
Each guard's active set is a down-set.
Thus \(g_d\) is nondecreasing in \(x,y\),
guard jumps included.

Let \(w_i\) be an outward upper bound for
\(\log(x_i/x_{i-1})\), on the \(2^{-128}\) lattice.
Let \(v_{ij}\) enclose \(g_{h_G}(x_i,x_j)\) upwards
on that lattice.
The diagonal cell has half the logarithmic square's area.
Hence
\begin{equation}\label{eq:app-upper-corners}
 J_{h_G}\le J_h^+
 =\sum_{i=1}^n w_i\left(
 \sum_{j<i}w_jv_{ij}+\frac12w_iv_{ii}\right).
\end{equation}

Split the inner interval at every affine crossing.
Include every guard crossing and each branch kink.
On each resulting interval the argument is \(A/z+B\).
If \(\Phi(t)=\int_1^t(1-f(s))\,ds\), then
\begin{equation}\label{eq:app-inner-primitive}
 \int_p^q\frac{1-f(A/z+B)}{z^2}\,dz
 =\frac{\Phi(A/p+B)-\Phi(A/q+B)}{A}
 \quad(A\ne0).
\end{equation}
The trivial bound is \(1/p-1/q\).
The smaller of the two upper bounds may be used.
For \(A=0\), the argument is constant.

The delay equations give the primitive explicitly:
\begin{equation}\label{eq:app-Phi}
 P(t)=\frac{tF(t)}{2e^\gamma},\qquad
 \Phi(t)=t-1-2e^\gamma(P(t+1)-1)\quad(t\ge0).
\end{equation}
For \(t\ge30\), monotonicity gives
\begin{equation}\label{eq:app-Phi-tail}
 \Phi(30)\le\Phi(t)
 \le\Phi(30)+(t-30)(1-f(30)).
\end{equation}
We bound the trapezoid error. If $g$ is either $F$ or $f$ and
$[a,a+h_0]\subset[1,30]$ is one mesh panel, then
\begin{equation}\label{eq:app-trapezoid-remainder}
 \int_a^{a+h_0}g(s)\,ds-
 {h_0\over2}\bigl(g(a)+g(a+h_0)\bigr)
 =-{h_0^3\over12}g''(\xi),
 \qquad a<\xi<a+h_0,
\end{equation}
with one-sided interpretations at the integer breakpoints. The mesh has
$h_0=1/65536$ and every breakpoint is a mesh node. We prove the derivative
bound independently of the quadrature table. Write $C=2e^\gamma$. The
elementary inequality $\gamma\le H_6-\log6$ and the first three positive
terms of the atanh series give
\[
 \log2>\frac{842}{1215},\qquad \log3>\frac{263}{240},\qquad
 \gamma<\frac{12853}{19440}<\frac23<\log2.
\]
Thus $0<C<4$ and $\log2<1$. The defining functions satisfy
\[
 F(s)=C/s\ (1\le s\le3),\quad
 f(s)=0\ (1\le s\le2),\quad f(s)=C\log(s-1)/s\ (2<s\le3),
\]
\[
 sF(s)=C+\int_1^{s-1}f(t)\,dt,\qquad
 sf(s)=\int_1^{s-1}F(t)\,dt\quad(s>3).
\]
Induction on unit intervals gives $0\le F,f\le C$. If the delayed
functions are in this range, the integrals have length $s-2$, so
$F(s)\le C(s-1)/s\le C$ and $f(s)\le C(s-2)/s\le C$.
For either companion pair $(g,h)=(F,f)$ or $(f,F)$, differentiating on
open unit intervals gives
\[
 g'(s)=\frac{h(s-1)-g(s)}s,\qquad
 g''(s)=\frac{h'(s-1)}s+
        \frac{2(g(s)-h(s-1))}{s^2}\quad(s>3).
\]
For $f$, the first-derivative delay identity holds already for $t>2$.
For $F$ on $2<t<3$, use $F'(t)=-C/t^2$. Hence
$|F'(t)|,|f'(t)|\le C/t$ for $t>2$. Since $g(s),h(s-1)\in[0,C]$, for $s>3$
\[
 |g''(s)|\le\frac{C}{s(s-1)}+\frac{2C}{s^2}
 \le\frac{7C}{18}<\frac{14}{9}<12.
\]
On the initial intervals,
\[
 F''(s)=2C/s^3\ (1<s<3),\quad f''(s)=0\ (1<s<2),
\]
\[
 f''(s)=C\left(\frac{2\log(s-1)}{s^3}
       -\frac1{s(s-1)^2}-\frac2{s^2(s-1)}\right)\ (2<s<3).
\]
The formulas imply $0\le F''\le8$ on $[1,2]$,
$0\le F''\le1$ on $[2,3]$, and $-4\le f''\le\log2<1$ on $[2,3]$.
Inside the last parentheses each negative term has magnitude at most
$1/2$. The positive term is at most $\log2/4$. Therefore
$|g''|\le8<12$ on all smooth panel interiors of $[1,30]$, with one-sided
limits at the integer breakpoints.

The archived outward intervals and the independently computed exact rational
intervals are recorded with the reproducibility materials. The derivative
check starts from the proved invariant and then rebuilds the delay table.
It checks every computed unit-interval enclosure against the archived coarse
intervals. It confirms that the computed global bound is $\le12$, the
recorded value:
\[
\begin{array}{c|cc}
\text{range}&F''&f''\\ \hline
[1,2]&[0,8]&[0,0]\\
{}[2,3]&[-8,8]&[-4,4]\\
{}[3,30]&[-12,12]&[-12,12]
\end{array}
\]
The computed bound is 8. Hence the absolute local remainder is at most
$8h_0^3/12=(2/3)h_0^3\le h_0^3$. The verifier checks this radius
against the $h_0^3$ panel allowance before it accepts the regenerated
table.
The M5 computation builds the delay table with \texttt{DEN=65536},
\texttt{S=128}, 160-bit Arb
interval arithmetic, and 1,900,545 generated nodes. The shifted computation
then evaluates the exact \(n_{\mathrm{mesh}}=512\) triangular
grid, with \(512\cdot513/2=131,328\) inner-grid evaluations. It stores
one aggregated row per outer index. Hence \texttt{SHIFT-N512-fresh.json} contains 512
rows for the 131,328 evaluations.
The M5 computation uses no Taylor model and no width-\(1/32\) step.
The historical Taylor certificate provides background information and is not an input to this calculation. The table and
shifted rows are regenerated before the directed inclusion check. The
verifier compares recorded JSON fields and exact row values. It also checks
SHA-256 hashes of the canonical compact-JSON integer arrays used by the
real-inclusion verifier. These canonical-array hashes differ from hashes
of regenerated JSON files. File formatting is not used as a certificate.
Directed summation gives Table~\ref{tab:app-m5}.
The positive last row proves the strict cap. The 512 aggregated rows give the exact 131328-term sum
used for \(J_h^+\). The directed real-inclusion result records the directed
tail and the cap check. The reproducibility report documents the execution
boundary.
The two normalization identities follow by expansion.
\end{proof}

The proposition concerns the scalar functional only.
The prime-sum comparison is required separately in Section~\ref{sec:retained-side}.

\subsection{Unshifted comparison (historical and not used)}\label{subsec:app-g24}

\paragraph{Status: historical and non-certified.}
This historical unshifted comparison records an earlier \(n=8192\) calculation for
context only. No current computation reproduces it. None of its values,
table, constants or comparison is an input to the main proof.
The comparison concerns the unshifted functional
$G_*:=\mathcal G_0$. It is distinct from the main-line coefficient bound
$G_0=6088440363/10^9$. The main K2 and L3 comparisons use the latter
via \cite[(4.73), (4.75), (4.79), (A.10a)]{PaperII}. They do not use
$U_0$. The cited bound $G_*\le G_0$ is the coefficient used in the main
argument. It does not assert the separate, sharper $G_*\le U_0$ comparison.
This appendix proves that comparison only conditionally, as an optional one.
Its historical delay table was constructed separately, but no such table or
producer is part of the current proof package.
The panel width is \(h_0=1/65536\).
At 192-bit precision, each trapezoid receives radius \(h_0^3\).
On smooth panel interiors, \(|F''|,|f''|<12\).
The panels meet the integer breakpoints.
This bounds each local truncation remainder. The endpoint enclosures
pass through the delay recursion in Arb ball arithmetic.
Accumulated input radii are thus kept before outward rounding.
All table endpoints are then rounded outwards to \(2^{-128}\).
The primitive and true tail use \eqref{eq:app-Phi-tail}.
This table gives a historical directed bound for \(R_{\mathrm{up}}\).
It is kept only to reproduce the earlier non-certified calculation.
The separate earlier Taylor models are not part of the M5 computation.
Let \(F_0^+\ge F(30)\), \(f_0^-\le f(30)\),
and \(H_0^+\ge H_{\log}\) be these directed endpoints.
Define
\begin{equation}\label{eq:app-unshifted-tail}
 R_0=\mathcal R(F_0^+,f_0^-,H_0^+)
 \ge R_{\mathrm{up}}.
\end{equation}
The earlier finite calculation used \(n=8192\) in
\eqref{eq:app-m5-mesh}.
The outer sum has \(n(n+1)/2=33558528\) nodes.
Write \(J_0^+\) for that upper-corner sum.
Table~\ref{tab:app-g24-rows} groups its eight consecutive row ranges.
Each row lists a lower and upper bound for the sum.
The diagonal half-weights are already included.

\begin{table}[ht]
\centering\small
\caption{Historical non-certified contributions to the unshifted sum.}
\label{tab:app-g24-rows}
\begin{tabular}{@{}rrrr@{}}
\hline
Rows & nodes & lower endpoint & upper endpoint\\
\hline
\(1\)--\(1024\)&524800&0.000002727813799350&0.000002727813799351\\
\(1025\)--\(2048\)&1573376&0.001226034866061399&0.001226034866061400\\
\(2049\)--\(3072\)&2621952&0.019651743682986805&0.019651743682986806\\
\(3073\)--\(4096\)&3670528&0.080548405705743750&0.080548405705743751\\
\(4097\)--\(5120\)&4719104&0.146720281588383935&0.146720281588383936\\
\(5121\)--\(6144\)&5767680&0.200161682376550457&0.200161682376550458\\
\(6145\)--\(7168\)&6816256&0.242188572133550572&0.242188572133550573\\
\(7169\)--\(8192\)&7864832&0.275672089024516687&0.275672089024516688\\
\hline
\end{tabular}
\end{table}

The earlier, non-certified calculation reported
\begin{equation}\label{eq:app-g24-numbers}
\begin{aligned}
 J_0^+&<0.966171537191592960294580903427,\\
 R_0&<0.000004373236943357822206406687,\\
 U_0&=1/\mu_*+R_0+J_0^++\frac1{2\cdot10^{85}}\\
 &\in[6.086175910428536318116787310112545656706640625,\\
 &\hspace{2.4em}6.086175910428536318116787310112545656706640626].
\end{aligned}
\end{equation}
The term \(1/(2\cdot10^{85})\) bounds directed outer-sum rounding.
The shifted computation does not use this interval.
The comparison with the continuous functional is one-sided:
\begin{equation}\label{eq:app-g24-inclusion}
 \text{Historical comparison (not certified): }G_* = \mathcal G_0\le U_0.
\end{equation}
No claim $2e^{-\gamma}U_0\le c$ follows from the printed interval.
The finite row sums alone prove no prime-sum estimate.
The source comparison and its error terms remain required.

\subsection{The sixth-sum constant}\label{subsec:app-roota}

The certified finite enclosure is
\begin{equation}\label{eq:app-roota-enclosure}
 \mathfrak g_{\mathrm{RA}}
 \in[0.0989598353353,\ 0.0989598353355].
\end{equation}
Here \(\mathfrak g_{\mathrm{RA}}\) denotes the integral quantity
in Section~\ref{sec:root-a}.
The constant used in the final bound is
\begin{equation}\label{eq:app-roota-floor}
 g_{\mathrm{RA,lo}}=
 \frac{989598353}{10^{10}}=0.0989598353.
\end{equation}
Thus the certified enclosure gives
\begin{equation}\label{eq:app-roota-comparison}
 \frac{7499}{9999}
 \left(J_J(1/2)-J_J(5227/10000)\right)
 -\frac{5106819}{31557955000}
 =\mathfrak g_{\mathrm{RA}}>g_{\mathrm{RA,lo}}.
\end{equation}
Only this conservative lower constant enters the final bound.

\subsection{Computer checks}\label{subsec:app-checks}

Table~\ref{tab:app-checks} lists the supplementary scripts.
The second column states what each one checks.
The partitions and defining formulas are printed above.
The exact-arithmetic checks use Python's integer and fraction types.
Environment and timing details are recorded in the reproducibility
materials. They are not theorem inputs.
A full numerical regeneration is distinguished from a compact check.
The certificate bundle is available at the location given in the
Data availability statement; its SHA-256 manifest hash is
\texttt{81fc00441c3a3f0924995cd41263feddece827d9b40376dae531191c86f90ed3}.
The recorded M5 regeneration agrees with the exact 512-row certificate.
The reproducibility certificate set is independent. It consists of the H6
\texttt{RESULT.json}, the 512-row M5 output, the Root--A enclosure, and the
verifier scripts for those records. Each record is checked against the source
hashes and the manifest. None replaces an analytic estimate.

\begin{center}\footnotesize
\begin{minipage}{\textwidth}\centering
\captionof{table}{Scripts and finite checks.}\label{tab:app-checks}
\setlength{\tabcolsep}{4pt}
\begin{tabular}{>{\raggedright\arraybackslash}p{.25\textwidth}>{\raggedright\arraybackslash}p{.54\textwidth}>{\raggedleft\arraybackslash}p{.09\textwidth}}
Script & Statements checked & Time\tabularnewline\hline
\noalign{\vskip 2pt}
\shortstack[l]{Reproducibility\\replay verifier}
& manifests, source hashes, TeX structure, finite consistency checks,
 and the producer verification run
 & replay\tabularnewline
\shortstack[l]{H6/M5 producer\\verification}
 & executes the H6 computation and M5 $n_{\mathrm{mesh}}=512$ table, shifted probe
 and real-inclusion verifier, computing the derivative enclosure and comparing
 regenerated outputs
 & replayed\tabularnewline
\shortstack[l]{H6 directed\\verifier}
 & regenerates the directed H6 rows from the source inputs and Arb runtime
 & replayed\tabularnewline
\shortstack[l]{M5 derivative\\verifier}
 & exact rational unit-interval bounds for $F''$ and $f''$ and the panel radius
 & replayed\tabularnewline
\shortstack[l]{M5 table, shift and\\inclusion verifiers}
 & regenerates the $DEN=65536$ delay table, $n_{\mathrm{mesh}}=512$ rows, and real inclusion
 & replayed\tabularnewline
\shortstack[l]{Finite consistency\\verifier}
 & exact finite arithmetic, orbit, and sign checks for the patches
 & replayed\tabularnewline
\end{tabular}
\end{minipage}
\end{center}

The reproducibility verifier checks the recorded inputs and then runs the
producer verification.
That run regenerates the H6 output, the M5 delay table, the shifted rows
and the real inclusion on the recorded platform.
The earlier unshifted calculation used Python~3.12.3
and python-flint~0.9.0, with four workers.
Its slowest worker took \(7639\) seconds, rounded upwards.
The historical unshifted comparison is optional and not certified.
The main proof does not use it, and it does not affect the shifted
computation.
Neither these scripts nor their agreement establishes a uniform estimate.
The final proof requires the analytic comparisons in Sections~\ref{sec:retained-side}--\ref{sec:l3}.

\subsection{Rows of the retained side}\label{subsec:app-retB-rows}

The index set $\mathcal I$ of Section~\ref{sec:retained-side} consists of
Paper~II's nineteen original row chains and one outer chain. The outer
chain $\mathrm{T5}$ covers the upper part of $S_5$, Lemma~5.1. Sixteen of the row
chains are listed below, each
with its source parent and raw coefficient. The remaining three are
Paper~II's single-prefix comparisons for $S_3$, $S_8$ and $S_9$. They keep
Paper~II's own sources, restorations and error entries unchanged.
The four rows of $\mathcal H$ are restored to their parents
$S_4$, $S_5$, $S_7$ and $S_{10}$. The lower part of $S_5$ has no separate
restoration address. The four joint rows are restored once to
the joint source of $G_1$ with weights $(2,1,1)$.
\begin{center}
\small
\begin{tabular}{@{}llr|llr@{}}
\hline
Row & Parent & Coefficient & Row & Parent & Coefficient\\
\hline
S1-0 & S1 & $3$ & S4 & S4 & $-1$\\
S1-1 & S1 & $-3$ & S5-low & S5 & $-1$\\
S1-2 & S1 & $3$ & S7 & S7 & $-1$\\
S1-3 & S1 & $-3$ & S10-Ret & S10 & $1$\\
S2-0 & S2 & $1$ & G1-0 & G1 $(2,1,1)$ & $-1$\\
S2-1 & S2 & $-1$ & G1-1 & G1 $(2,1,1)$ & $1$\\
S2-2 & S2 & $1$ & G1-2 & G1 $(2,1,1)$ & $-1$\\
S2-3 & S2 & $-1$ & G1-3 & G1 $(2,1,1)$ & $1$\\
\hline
\end{tabular}
\end{center}
These sixteen rows, their seven source-parent restorations and four outer
terms carry $91$ error addresses. Each address occurs exactly once. One
outer term is the chain $\mathrm{T5}$. The others are the strict BV piece of
$S_2$, inside its joint group, and two $G_1$ comparator entries. The
same table has seven further rows with $38$ addresses. They treat
$S_3$, $S_8$ and $S_9$ another way, and this paper does not use them.
\subsection{Statements for the sixth sum}\label{subsec:app-ra-rows}

The sixth-sum argument uses the following statements.
\begin{center}
\small
\begin{tabular}{@{}>{\raggedright\arraybackslash}p{0.96\textwidth}@{}}
\hline
Statement\\
\hline
Lemma~6.2 (rough-factor distribution)\\
Lemma~6.4 (four-family Siegel--Walfisz); range \(N^{1/3}/8\le H\)\\
Lemmas~6.6 and~6.7 (selector and tensor, applied per shell pair)\\
Lemma~6.5 (transferred mean square)\\
Lemma~6.8 (two-factor convolution)\\
Lemma~6.9 (central rectangles)\\
Proper powers, (6.26)\\
Fact~3.2 (maximal prime estimates)\\
Paper~I, (2.21), as used in Paper~I, (3.25)\\
\hline
\end{tabular}
\end{center}
